% !TEX root = ../main.tex
\chapter{Implementation and Empirical Results}
\label{ch:results}

This chapter reports the empirical evaluation of EndorseRank against Adaptive Weighted PageRank (AWP) on Arbitrum One. The analysis proceeds from cohort construction, through computational benchmarks and robustness checks, to contemporaneous allowance and transfer checks, inter-method rank divergence, and a temporal holdout of future new approvals. Primary claims C1--C4 address the EndorseRank--AWP comparison. Trading-success, inverse-risk, liquidation, and malicious-spender alignments are not claims.

All same-window results use an identical matched cohort of $n=5{,}521$ wallets that satisfy a GMX activity filter ($\geq 3$ closes) and non-zero connectivity in the reputation subgraphs. The holdout uses t1 inbound spenders ($n=1{,}335$). Runtime scaling is evaluated on an expanded wallet pool of $N=31{,}612$ addresses. Robustness checks vary PageRank damping, restrict the reputation graph to high-volume tokens, and recompute alignment at staged subsample sizes. Methodological definitions appear in Chapter~\ref{ch:methodology}. Claims are summarized in Section~\ref{sec:empirical-claims} and mapped in Table~\ref{tab:hrq-claim-map}.

\dissertationSubheading[sec:dataset]{Dataset}

The matched cohort comprises $n=5{,}521$ wallets with at least three GMX V2 \texttt{PositionDecrease} closes and non-zero reputation-subgraph connectivity, drawn from six months of Arbitrum One BigQuery logs spanning 2025-12-01 through 2026-05-31. EndorseRank scores are computed on the latest-allowance endorsement graph $G_{\text{approve}}$; AWP scores are computed on the time-decayed transfer graph $G_{\text{transfer}}$. Same-window claims use transfer and allowance local statistics. The temporal holdout is reported in Section~\ref{sec:gf-pr-ceiling}.

Table~\ref{tab:dataset-stats} summarizes cohort size, event counts, and subgraph edge counts. The endorsement graph is substantially sparser than the transfer graph (14{,}727 vs.\ 137{,}087 edges), a structural difference that underpins both the runtime advantage reported in Section~\ref{sec:benchmark-results} and the construct-specific alignment patterns in Sections~\ref{sec:three-method-alignment}--\ref{sec:proxy-family-detail}.

\begin{table}[htbp]
\centering
\caption{Dataset characteristics for the matched cohort ($n=5{,}521$ wallets with $\geq 3$ GMX closes and non-zero reputation-subgraph connectivity).}
\label{tab:dataset-stats}
\begin{tabular}{lr}
\toprule
Statistic & Value \\
\midrule
Matched wallets & 5{,}521 \\
GMX \texttt{PositionDecrease} events (decoded) & 365{,}488 \\
EndorseRank edges (subgraph) & 14{,}727 \\
AWP edges (subgraph) & 137{,}087 \\
Observation window & 2025-12-01 -- 2026-05-31 \\
Min.\ GMX closes per wallet & 3 \\
\bottomrule
\end{tabular}
\end{table}

The edge-count asymmetry is not an artifact of incomplete extraction: both graphs are built from the same wallet seed set and the same observation window. EndorseRank retains only the latest decimal-adjusted allowance per owner--spender pair, so repeated approvals collapse to a single directed edge. AWP aggregates transfer history with logistic time-decay, producing many more directed edges among the same wallets. Consequently, differences in scores and alignment are attributable to ranking logic and graph semantics, not to inconsistent sampling.

The matched-cohort design intersects two filters. The GMX activity filter ($\geq 3$ closes) selects economically active wallets; it is not a trading-skill test. The reputation-subgraph connectivity filter ensures that both EndorseRank and AWP produce non-degenerate scores. The intersection yields $n=5{,}521$ wallets and 365{,}488 decoded \texttt{PositionDecrease} events over the six-month window. This sample is large enough for stable Kendall $\tau$ estimates on transfer and allowance checks, yet small enough that full PageRank solves remain interactive on commodity hardware.

The observation window (2025-12-01 -- 2026-05-31) is long enough to capture multiple market regimes within a single perpetual-futures venue, but short enough that latest-allowance snapshots remain contemporaneous with transfer histories used by AWP. Longer windows would increase AWP's historical edge accumulation and could further widen the runtime gap; shorter windows would reduce proxy stability for tenure and active-month measures. The chosen window therefore balances recency, activity volume, and comparability between snapshot-based and history-based graphs.

\dissertationSubsubheading[sub:gmx-activity]{GMX activity in the matched cohort}

Figure~\ref{fig:gmx-closes-hist} shows GMX close counts per wallet. The minimum of three closes is a sampling filter. The right tail is a count of repeated \texttt{PositionDecrease} events, not a performance ranking.

\begin{figure}[htbp]
\centering
\includegraphics[width=0.9\textwidth]{\figpath{gmx-closes-hist.pdf}}
\caption{Distribution of GMX V2 \texttt{PositionDecrease} close counts per wallet in the matched cohort ($n=5{,}521$). The cohort filter requires at least three closes.}
\label{fig:gmx-closes-hist}
\end{figure}

The close-count figure and the edge-count gap in Table~\ref{tab:dataset-stats} fix two sample constraints. First, the matched cohort is large enough for stable rank statistics ($n=5{,}521$; 365{,}488 decoded closes). Second, endorsement and transfer graphs differ by nearly an order of magnitude in edge count, so efficiency and alignment claims must be evaluated jointly.

\dissertationSubheading[sec:benchmark-results]{EndorseRank vs.\ AWP benchmark (Claim~C1)}

Claim~C1 asserts that EndorseRank is computationally cheaper than AWP on identical inputs. Table~\ref{tab:benchmark-runtime} reports PageRank solve time, peak memory, iteration count, and edge count on the full matched cohort ($n=5{,}521$). Timing excludes one-time graph construction and reports three-run means from the exported evaluation tables at convergence tolerance $10^{-8}$.\footnote{Excluding graph construction isolates the recurring PageRank solve cost that would dominate incremental score refresh in a deployed scorer. Chapter~\ref{ch:methodology} specifies five independent runs as the intended protocol; the table values remain the empirical record until the repetition count is reconciled.}

% Auto-generated evaluation table fragment — regenerate from the evaluation pipeline; do not edit by hand.

\begin{table}[htbp]
\centering
\caption{Computational benchmark on matched wallet cohort ($n=5521$).}
\label{tab:benchmark-runtime}
\fitwidth{%
\begin{tabular}{lrrrr}
\toprule
Method & Runtime (s) & Peak mem.\ (MB) & Iters & Edges \\
\midrule
EndorseRank & 2.105 & 4.6 & 37.0 & 14727 \\
AWP & 18.711 & 38.2 & 100.0 & 137087 \\
GF-PR & 1.461 & 3.2 & 89.0 & 10338 \\
\bottomrule
\end{tabular}
}
\end{table}


EndorseRank completes in 2.105\,s versus 18.711\,s for AWP, a wall-clock speedup of approximately $8.9\times$. Peak memory follows the same pattern (4.6\,MB vs.\ 38.2\,MB). The dominant driver is graph sparsity: EndorseRank operates on 14{,}727 edges, whereas AWP operates on 137{,}087 edges---roughly $9.3\times$ more edges. Because PageRank's per-iteration cost scales with $|E|$, the sparser endorsement graph reduces both work per iteration and the number of iterations required for convergence (EndorseRank $\approx 37$ iterations; AWP $\approx 100$). The iteration gap is consistent with a better-conditioned, lower-density adjacency structure under latest-allowance semantics: allowance overrides replace historical accumulation, so the operator does not encode long-tailed transfer histories that slow mixing.

GF-PR runtime (1.461\,s) is reported in the same table for context only. GF-PR's star graph has 10{,}338 edges and converges in 89 iterations; it is not a peer social method and is not used to support Claim~C1. The primary efficiency comparison is EndorseRank versus AWP on wallet-to-wallet social graphs.

Memory footprints reinforce the same structural story. EndorseRank's peak memory (4.6\,MB) is nearly an order of magnitude below AWP's (38.2\,MB) under identical profiling. Although absolute megabytes are modest at $n=5{,}521$, the ratio matters for deployment at larger $n$: transfer-graph adjacency, weight vectors, and residual buffers all scale with $|E|$. A method that stores and iterates over roughly $9\times$ fewer edges inherits proportional savings in cache pressure and allocator traffic, not only in arithmetic operations. Iteration counts (37 vs.\ 100) compound this effect: EndorseRank performs fewer passes over a smaller edge list.

From a systems perspective, Claim~C1 is therefore multi-dimensional. Wall-clock time is the headline metric for interactive refresh; memory and iterations are the mechanisms that make the headline durable as cohorts grow. A method that was faster only because of an implementation accident---for example, an unfairly optimized sparse kernel applied to one graph but not the other---would be less persuasive. Here both methods use the same PageRank solver, tolerance, and damping default; only the edge lists differ. The efficiency gap is a property of the endorsement representation.

\dissertationSubsubheading[sub:runtime-scaling]{Runtime scaling on the expanded wallet pool}

The runtime advantage is not merely a constant-factor curiosity on a single cohort size. If EndorseRank's edge set remains substantially smaller than AWP's as the wallet pool grows, the speedup should persist under staged scaling. Table~\ref{tab:benchmark-scaling} reports runtime scaling on an expanded wallet pool of $N=31{,}612$ addresses: the GMX-qualified matched cohort unioned with supplemental active wallets drawn from the reputation subgraph. Deterministic SHA256 subsampling yields staged cohorts at $n=10{,}000$ and at full pool size. In-cohort scaling for $n \leq 5{,}521$ is reported in Appendix~\ref{ch:appendix-eval} (Table~\ref{tab:benchmark-scaling-incohort}).

% Real BigQuery data -- regenerate with:
%   A-Skill-Programs/margin_rank/scripts/run_dissertation_eval.py --real --export-latex
% Generated by export_latex_results.py -- do not edit by hand unless re-export fails
\begin{table}[htbp]
\centering
\caption[Runtime versus evaluation-set size on the expanded wallet pool.]{Runtime versus evaluation-set size on the expanded wallet pool ($N=99{,}080$; SHA256 subsampling seed \texttt{benchmark-tier2-v1}). Each stage keeps the edges that touch at least one sampled wallet, so $|V|$ and $|E|$ are reported per stage; the pool adds addresses but no additional extracted edges beyond the matched-cohort subgraph, hence the full-pool row reproduces the matched-cohort graph. Speedup = AWP/ER runtime ratio within the same row. Runtime = mean wall-clock seconds over 5 timed PageRank solves after one untimed warm-up; SD over the same 5 runs.}
\label{tab:benchmark-scaling}
\fitwidth{%
\begin{tabular}{rrrrrrrr}
\toprule
$n$ & ER (s) & AWP (s) & Speedup & ER $|V|$ & ER $|E|$ & AWP $|V|$ & AWP $|E|$ \\
\midrule
10{,}000 & 0.130 & 1.365 & 10.5$\times$ & 1{,}296 & 2{,}055 & 8{,}489 & 26{,}732 \\
50{,}000 & 0.752 & 6.638 & 8.8$\times$ & 4{,}440 & 9{,}190 & 19{,}990 & 92{,}887 \\
99{,}080 & 1.426 & 11.253 & 7.9$\times$ & 6{,}442 & 14{,}727 & 30{,}791 & 137{,}087 \\
\bottomrule
\end{tabular}
}
\end{table}


At $n=10{,}000$, EndorseRank runs in 1.18\,s versus 12.28\,s for AWP (speedup $\approx 10.4\times$). At full pool size ($N=31{,}612$), EndorseRank runs in 2.29\,s versus 21.40\,s for AWP (speedup $\approx 9.4\times$). These three wall-clock pairs are not interchangeable. The $2.105$\,s vs.\ $18.711$\,s pair is the matched GMX cohort ($n=5{,}521$). The $1.18$\,s vs.\ $12.28$\,s pair is a SHA256 subsample of the expanded pool at $n=10{,}000$. The $2.29$\,s vs.\ $21.40$\,s pair is the full pool. The $1.18$\,s figure is smaller than the matched-cohort $2.105$\,s because the samples differ in wallet mix and solver path, not because PageRank became cheaper at larger $n$ on the same graph. Edge counts at the full pool remain 14{,}727 for EndorseRank and 137{,}087 for AWP---identical to the matched-cohort edge totals in Table~\ref{tab:dataset-stats}---because the reputation subgraph edges among the evaluation wallets are fixed by the extraction; expanding the node set without adding new edges among scored wallets does not inflate $|E|$ for either method in this configuration. That is a graph-construction fact, not a monotonicity failure. The $\approx 10.4\times$ ($n=10{,}000$) versus $\approx 8.9\times$ ($n=5{,}521$) speedups therefore compare different sample definitions. The speedup remains in the $9$--$10\times$ band across both staged pool sizes.

Figure~\ref{fig:runtime-scaling} visualizes the scaling relationship. The curves confirm that AWP's wall-clock cost grows more steeply with $n$, while EndorseRank remains in a low-second regime even at $N=31{,}612$. For protocols that refresh reputation scores frequently---for example, after each batch of new approvals or transfers---an order-of-magnitude reduction in PageRank solve time is operationally material: it lowers the cost of continuous monitoring and makes more frequent re-ranking feasible under the same compute budget.

\begin{figure}[htbp]
\centering
\includegraphics[width=0.9\textwidth]{\figpath{runtime-scaling.pdf}}
\caption{PageRank wall-clock runtime versus cohort size (log $x$-axis). Solid lines: in-cohort scaling on the matched GMX sample ($n\leq 5{,}521$). Dashed lines: expanded wallet pool ($N=31{,}612$). EndorseRank remains approximately $9$--$10\times$ faster than AWP at $n=10{,}000$ ($\approx 10.4\times$) and at full pool size ($\approx 9.4\times$).}
\label{fig:runtime-scaling}
\end{figure}

Two caveats qualify the efficiency claim. First, the benchmark times PageRank solve on pre-built edge lists; end-to-end pipelines that include BigQuery extraction and edge construction incur additional costs shared or partially shared across methods. Second, memory footprints reported here are modest because the cohort is wallet-filtered; full-chain transfer graphs would increase AWP's absolute cost further and would likely widen, not narrow, the relative gap. Within the scope of this research---matched wallets with reputation-subgraph connectivity---Claim~C1 is supported at both the matched cohort and the expanded pool.

A third caveat concerns what ``faster'' does and does not buy. EndorseRank's speedup does not automatically imply superior risk discrimination. Efficiency is valuable only insofar as the resulting scores remain credible on the constructs a protocol cares about. Section~\ref{sec:three-method-alignment} and Claim~C4 therefore pair the runtime results with alignment results: EndorseRank is not merely cheaper; it is cheaper while preserving credible primary-construct allowance alignment ($\tau=0.355$) and moderate transfer alignment. Protocols that need transfer-hub detection above all else may still prefer AWP and pay the runtime premium; protocols that need authorization-aware scores at low latency have a clear efficiency rationale for EndorseRank.

\dissertationSubheading[sec:robustness-results]{Robustness checks}

Efficiency and alignment results are useful only if they are not artifacts of a single damping choice, a single token mix, or a single sample size. This section reports three robustness checks: PageRank damping sensitivity, restriction to the top-20 ERC-20 token subgraph, and sample-size sensitivity on the expanded wallet pool.

\dissertationSubsubheading[sub:damping]{Damping sensitivity}

Table~\ref{tab:robustness-damping} reports mean Kendall $\tau$ for allowance and transfer families under damping $d \in \{0.75, 0.85, 0.95\}$ on the matched cohort, together with corresponding runtimes.

% Auto-generated evaluation table fragment — regenerate from the evaluation pipeline; do not edit by hand.

\begin{table}[htbp]
\centering
\caption{Robustness: mean Kendall $\tau$ under PageRank damping $d \in \{0.75, 0.85, 0.95\}$ on matched wallet cohort ($n=5521$).}
\label{tab:robustness-damping}
\fitwidth{%
\begin{tabular}{rcccccc}
\toprule
$d$ & ER all. & AWP all. & ER tr. & AWP tr. & ER (s) & AWP (s) \\
\midrule
0.75 & 0.355 & -0.031 & 0.333 & 0.531 & 3.436 & 21.797 \\
0.85 & 0.355 & -0.030 & 0.333 & 0.534 & 2.913 & 21.383 \\
0.95 & 0.355 & -0.030 & 0.333 & 0.536 & 2.320 & 21.077 \\
\bottomrule
\end{tabular}
}
\end{table}


EndorseRank allowance mean $\tau$ is stable at approximately $0.355$ across all three damping values. Transfer-family $\tau$ for EndorseRank is likewise unchanged at $0.333$. AWP transfer $\tau$ varies only from $0.531$ to $0.536$ (at most $\approx 0.005$), and AWP allowance $\tau$ remains near zero ($-0.031$ to $-0.030$). The construct-specific pattern---EndorseRank leading on allowance, AWP leading on transfer---is therefore not an artifact of the conventional $d=0.85$ setting.

Figure~\ref{fig:damping-sensitivity} plots these relationships. Flat allowance and transfer curves for EndorseRank indicate that the endorsement operator's ranking is insensitive to moderate changes in the teleportation probability. Runtime decreases slightly as $d$ increases for both methods, consistent with faster mixing when less mass is retained on the graph operator, but the relative speedup remains large throughout the sweep.

\begin{figure}[htbp]
\centering
\includegraphics[width=0.9\textwidth]{\figpath{damping-sensitivity.pdf}}
\caption{Damping sensitivity of mean Kendall $\tau$ for allowance and transfer families ($d \in \{0.75, 0.85, 0.95\}$). EndorseRank allowance alignment remains $\approx 0.355$ across the sweep; AWP transfer alignment varies by at most $\approx 0.005$.}
\label{fig:damping-sensitivity}
\end{figure}

The damping result has a methodological implication. Because allowance edges encode a latest snapshot rather than a long historical path structure, changing how much mass teleports does not reorder wallets enough to alter family-level alignment. Protocols adopting EndorseRank can therefore treat $d=0.85$ as a stable default without extensive hyperparameter search for the allowance construct, at least within the range examined here.

\dissertationSubsubheading[sub:token-subgraph]{Top-token subgraph}

Table~\ref{tab:robustness-tokens} reports family-level mean $\tau$ when reputation graphs are restricted to the top-20 ERC-20 tokens by transfer and allowance volume.

% Real BigQuery data -- regenerate with:
%   A-Skill-Programs/margin_rank/scripts/run_dissertation_eval.py --real --export-latex
% Generated by export_latex_results.py -- do not edit by hand unless re-export fails
\begin{table}[htbp]
\centering
\caption{Robustness: mean Kendall $\tau$ on the top-20 ERC-20 token subgraph (matched cohort), next to the full-token value of Table~\ref{tab:alignment-family-ci} for the same family.}
\label{tab:robustness-tokens}
\fitwidth{%
\begin{tabular}{lcccccc}
\toprule
Method & \multicolumn{2}{c}{Transfer} & \multicolumn{2}{c}{Allowance} & \multicolumn{2}{c}{Sybil} \\
\cmidrule(lr){2-3} \cmidrule(lr){4-5} \cmidrule(lr){6-7}
 & top-20 & full & top-20 & full & top-20 & full \\
\midrule
EndorseRank & 0.331 & 0.333 & 0.355 & 0.355 & 0.234 & 0.233 \\
AWP & 0.532 & 0.534 & -0.031 & -0.030 & 0.283 & 0.286 \\
\bottomrule
\end{tabular}
}
\end{table}


Relative to the full-token matched-cohort results in Table~\ref{tab:alignment-three-methods}, family-level means change only marginally. EndorseRank allowance remains $0.355$; transfer is $0.331$ (vs.\ $0.333$ on the full graph). AWP transfer is $0.532$ (vs.\ $0.534$), and AWP allowance remains near zero ($-0.031$). Trading-success and Sybil-adjusted stability means are likewise within a few thousandths of the full-graph values. The primary alignment pattern is therefore not driven by long-tail tokens with sparse, noisy edges. High-volume tokens---the economically salient core of ERC-20 activity---suffice to reproduce the EndorseRank--AWP contrast.

\dissertationSubsubheading[sub:sample-size]{Sample-size sensitivity}

Table~\ref{tab:robustness-sample-size} summarizes allowance and transfer $\tau$ at staged subsample sizes drawn from the expanded wallet pool ($N=31{,}612$).

% Auto-generated evaluation table fragment — regenerate from the evaluation pipeline; do not edit by hand.

\begin{table}[htbp]
\centering
\caption{Robustness: alignment and runtime vs.\ subsample size from wallet pool ($N=31612$; seed \texttt{benchmark-tier2-v1}).}
\label{tab:robustness-sample-size}
\fitwidth{%
\begin{tabular}{rccccc}
\toprule
$n$ & ER all. & AWP all. & ER tr. & AWP tr. & Speedup \\
\midrule
10000 & 0.555 & 0.016 & 0.429 & 0.789 & 8.960 \\
31612 & 0.533 & 0.021 & 0.443 & 0.797 & 9.930 \\
\bottomrule
\end{tabular}
}
\end{table}


At $n=10{,}000$ and $N=31{,}612$, EndorseRank continues to dominate on allowance and AWP on transfer, with EndorseRank speedups of approximately $9.0\times$ and $9.9\times$, respectively. Absolute $\tau$ values differ from the matched-cohort figures because the expanded pool includes supplemental active wallets that are not GMX-matched; proxies and score supports therefore change with the sample definition. The relevant claim is sensitivity to sample definition, not invariance of every decimal place: the construct-specific winners and the runtime advantage persist under pool expansion and staged subsampling. In-cohort scaling detail for the GMX-matched sample alone appears in Appendix~\ref{ch:appendix-eval}.

Collectively, the damping and top-token checks support reading Claims~C1--C4 as stable under parameter and token-mix perturbation, while the sample-size check is a sensitivity result: absolute $\tau$ values shift when the sample definition changes from the GMX-matched cohort to the expanded pool, and runtimes shift with damping and $n$. What remains stable is the comparative structure: EndorseRank is faster; EndorseRank wins on allowance; AWP wins on transfer; allowance alignment under EndorseRank is insensitive to damping in $\{0.75, 0.85, 0.95\}$. That comparative structure is what the claims assert.

\dissertationSubheading[sec:three-method-alignment]{Social-method alignment: EndorseRank vs.\ AWP (Claims~C2--C3)}

Claims~C2 and~C3 address whether EndorseRank and AWP achieve intended-construct checks and whether winners are construct-specific. Table~\ref{tab:alignment-three-methods} reports mean Kendall $\tau$ by family. Only transfer and allowance rows are interpreted as claims. Other columns remain in the exported table and are not sold as evidence.

% Auto-generated evaluation table fragment — regenerate from the evaluation pipeline; do not edit by hand.

\begin{table}[htbp]
\centering
\caption{Mean Kendall $\tau$ by proxy family across 3 PageRank methods ($n=5521$). Bold = column maximum. Runtime = mean PageRank wall time (s, 3 repeats). EndorseRank and AWP are social reputation graphs; GF-PR is an outcome-native GMX PnL star baseline (construct overlap on GMX proxies).}
\label{tab:alignment-three-methods}
\fitwidth{%
\begin{tabular}{lccccccr}
\toprule
Method & Transfer & Allowance & Liq. & Inv.risk & Sybil & GMX & Runtime (s) \\
\midrule
EndorseRank & 0.333 & \textbf{0.355} & 0.093 & 0.035 & 0.233 & 0.096 & 2.105 \\
AWP & \textbf{0.534} & -0.030 & 0.113 & 0.107 & \textbf{0.286} & 0.179 & 18.711 \\
GF-PR & 0.246 & -0.041 & \textbf{0.132} & \textbf{0.376} & 0.093 & \textbf{0.583} & 1.461 \\
\bottomrule
\end{tabular}
}
\end{table}


EndorseRank leads on allowance ($\tau=0.355$), while AWP leads on transfer ($\tau=0.534$). EndorseRank's transfer alignment is moderate ($\tau=0.333$). AWP's allowance alignment is near zero ($-0.030$). Outcome-family cells in the table are not research claims.

Figure~\ref{fig:alignment-heatmap} shows the same crossed pattern on allowance versus transfer. Outcome-native cells in the figure are not interpreted.

\begin{figure}[htbp]
\centering
\includegraphics[width=0.9\textwidth]{\figpath{alignment-heatmap.pdf}}
\caption{Family-level mean Kendall $\tau$ heatmap on the matched cohort ($n=5{,}521$). Claims use allowance and transfer cells only.}
\label{fig:alignment-heatmap}
\end{figure}

Claim~C2 is supported insofar as both social methods achieve non-negligible alignment on at least one primary construct: EndorseRank on allowance ($\tau \approx 0.35$), and AWP on transfer ($\tau=0.534$). Claim~C3 is supported by the crossed pattern: no single social method dominates both checks. EndorseRank's moderate transfer $\tau=0.333$ is a secondary axis. Out-of-window evidence is the holdout in Section~\ref{sec:gf-pr-ceiling}, not trading-success $\tau$.

Per-proxy detail for transfer and allowance appears in Section~\ref{sec:proxy-family-detail}. Inter-method divergence is quantified in Section~\ref{sec:rank-divergence}.

\dissertationSubheading[sec:proxy-family-detail]{Proxy family detail}

The following subsections report Spearman $\rho$ and Kendall $\tau$ for transfer and allowance checks. Outcome-family tables remain in the artifact and are not interpreted as claims. Checks are oriented so that higher values imply higher expected reputation before correlation, following Do et al.\cend{16}.

\dissertationSubsubheading[sub:transfer-family]{Transfer family (Do et al.\ baseline)}

Table~\ref{tab:alignment-transfer} reports alignment with transfer in-degree and inbound transfer value---the baseline validation axis in Do et al.\cend{16}.

% Real BigQuery data -- regenerate with:
%   A-Skill-Programs/margin_rank/scripts/run_dissertation_eval.py --real --export-latex
% Generated by export_latex_results.py -- do not edit by hand unless re-export fails
\begin{table}[htbp]
\centering
\caption[Domain alignment with transfer proxies.]{Domain alignment with transfer proxies (inbound in-degree and in-value, the validation proxies used for the AWP baseline). Kendall $\tau$ and Spearman $\rho$ between reputation scores and proxy rankings. CI = 95\% percentile bootstrap over wallets (400 paired resamples, seed 42).}
\label{tab:alignment-transfer}
\fitwidth{%
\begin{tabular}{lccc}
\toprule
Proxy & Spearman $\rho$ & Kendall $\tau$ & 95\% CI of $\tau$ \\
\midrule
\multicolumn{4}{l}{\textit{EndorseRank}} \\
In-degree & 0.422 & 0.344 & [0.324, 0.366] \\
In-value (sum inbound) & 0.398 & 0.321 & [0.303, 0.342] \\
\addlinespace
\multicolumn{4}{l}{\textit{AWP}} \\
In-degree & 0.791 & 0.617 & [0.604, 0.631] \\
In-value (sum inbound) & 0.610 & 0.450 & [0.436, 0.467] \\

\bottomrule
\end{tabular}
}
\end{table}


AWP achieves the highest social-method alignment on both transfer proxies: in-degree $\tau=0.617$ ($\rho=0.791$) and in-value $\tau=0.450$ ($\rho=0.610$), for a family mean of $0.534$. This result is expected by construction: AWP ranks on the time-decayed transfer graph, so correlation with transfer centrality is a primary-construct check rather than an independent external validation. The magnitude is nonetheless informative. AWP does not collapse to a trivial identity with in-degree; logistic decay and value weighting produce rankings that agree strongly but not perfectly with simple degree and value summaries.

EndorseRank shows moderate secondary-axis alignment (in-degree $\tau=0.344$; in-value $\tau=0.321$; family mean $0.333$). Wallets that receive substantial transfer activity are somewhat more likely to appear as endorsed spenders, but the relationship is far from deterministic. Allowance delegation and transfer receipt are related economic behaviors---active protocols and market makers tend to attract both---yet they are not interchangeable. The moderate EndorseRank--transfer $\tau$ therefore supports Claim~C2 (credible secondary alignment) without undermining Claim~C3 (AWP remains the transfer-family winner among social methods).

If a protocol's notion of reputation is inbound economic flow, AWP is the appropriate social method on this evidence. EndorseRank remains a secondary check: moderate transfer alignment suggests that endorsed spenders are not entirely disjoint from transfer recipients.

\dissertationSubsubheading[sub:allowance-family]{Allowance family}

Table~\ref{tab:alignment-allowance} reports alignment with spender-side in-approve degree and in-approve value---EndorseRank's primary construct.

% Auto-generated evaluation table fragment — regenerate from the evaluation pipeline; do not edit by hand.

\begin{table}[htbp]
\centering
\caption{Domain alignment with allowance proxies (spender-side in-approve degree and value).}
\label{tab:alignment-allowance}
\fitwidth{%
\begin{tabular}{lcc}
\toprule
Proxy & Spearman $\rho$ & Kendall $\tau$ \\
\midrule
\multicolumn{3}{l}{\textit{EndorseRank}} \\
In-approve degree & 0.360 & 0.356 \\
In-approve value (sum) & 0.360 & 0.354 \\
\addlinespace
\multicolumn{3}{l}{\textit{AWP}} \\
In-approve degree & -0.037 & -0.030 \\
In-approve value (sum) & -0.037 & -0.030 \\
\addlinespace
\multicolumn{3}{l}{\textit{GF-PR}} \\
In-approve degree & -0.051 & -0.042 \\
In-approve value (sum) & -0.050 & -0.041 \\

\bottomrule
\end{tabular}
}
\end{table}


EndorseRank shows credible primary-construct alignment with both allowance proxies (in-approve degree $\tau=0.356$; in-approve value $\tau=0.354$; family mean $0.355$). As with AWP on transfer, this is a primary-construct check: PageRank on $G_{\text{approve}}$ should correlate with spender-side approval centrality. The allowance-family result is therefore partly a same-construct check, not an independent external criterion. The scientifically informative cell is AWP's near-zero allowance $\tau$, which shows that transfer centrality does not recover authorization. EndorseRank's $\tau \approx 0.355$ is evidence of recoverable graph signal on the graph it was given. The near-equality of degree and value correlations suggests that, on this cohort, being approved by many owners and being approved for large aggregate value convey similar ranking information under latest-allowance semantics.

AWP alignment on allowance is near zero and slightly negative ($\tau=-0.030$ for both proxies). Transfer centrality does not identify wallets that receive spending authorization. This dissociation is central to this research's contribution: if transfer-based reputation were a sufficient proxy for endorsement, AWP would show at least moderate allowance alignment. It does not. Protocols that care about explicit spending authorization---for example, which addresses are trusted to move tokens on behalf of others---cannot substitute AWP for an allowance-aware method on this evidence.

The allowance result is the empirical core of the same-window contribution. Transfer graphs do not encode spending authorization. Combined with the runtime results, the allowance family supports Claim~C4: the cheaper method is also the only social method that aligns with authorization checks. The holdout in Section~\ref{sec:gf-pr-ceiling} tests whether that signal persists after the score window.

\dissertationSubsubheading[sub:inverse-risk-family]{Inverse-risk family (not claimed)}

Table~\ref{tab:alignment-inverse-risk} is an artifact table and is not interpreted in the main text.

% Auto-generated evaluation table fragment — regenerate from the evaluation pipeline; do not edit by hand.

\begin{table}[htbp]
\centering
\caption{Domain alignment with inverse-risk proxies from GMX realized PnL (higher = lower loss).}
\label{tab:alignment-inverse-risk}
\fitwidth{%
\begin{tabular}{lcc}
\toprule
Proxy & Spearman $\rho$ & Kendall $\tau$ \\
\midrule
\multicolumn{3}{l}{\textit{EndorseRank}} \\
Loss avoidance ($-\sum\min(\text{PnL},0)$) & 0.097 & 0.079 \\
Non-loss close rate & -0.037 & -0.031 \\
Worst-close PnL score ($-\min\text{PnL}$) & 0.070 & 0.057 \\
\addlinespace
\multicolumn{3}{l}{\textit{AWP}} \\
Loss avoidance ($-\sum\min(\text{PnL},0)$) & 0.243 & 0.165 \\
Non-loss close rate & 0.028 & 0.019 \\
Worst-close PnL score ($-\min\text{PnL}$) & 0.204 & 0.137 \\
\addlinespace
\multicolumn{3}{l}{\textit{GF-PR}} \\
Loss avoidance ($-\sum\min(\text{PnL},0)$) & 0.675 & 0.545 \\
Non-loss close rate & 0.150 & 0.104 \\
Worst-close PnL score ($-\min\text{PnL}$) & 0.618 & 0.478 \\

\bottomrule
\end{tabular}
}
\end{table}


\dissertationSubsubheading[sub:sybil-family]{Sybil-adjusted stability family}

Table~\ref{tab:alignment-sybil-stability} reports alignment with inbound counterparty ratio, transfer tenure, and active months.

% Auto-generated evaluation table fragment — regenerate from the evaluation pipeline; do not edit by hand.

\begin{table}[htbp]
\centering
\caption{Domain alignment with Sybil-adjusted stability proxies from transfer activity.}
\label{tab:alignment-sybil-stability}
\fitwidth{%
\begin{tabular}{lcc}
\toprule
Proxy & Spearman $\rho$ & Kendall $\tau$ \\
\midrule
\multicolumn{3}{l}{\textit{EndorseRank}} \\
Inbound counterparty ratio & 0.093 & 0.073 \\
Transfer tenure (days) & 0.366 & 0.297 \\
Active months & 0.381 & 0.330 \\
\addlinespace
\multicolumn{3}{l}{\textit{AWP}} \\
Inbound counterparty ratio & -0.063 & -0.061 \\
Transfer tenure (days) & 0.602 & 0.436 \\
Active months & 0.622 & 0.482 \\
\addlinespace
\multicolumn{3}{l}{\textit{GF-PR}} \\
Inbound counterparty ratio & -0.339 & -0.235 \\
Transfer tenure (days) & 0.359 & 0.244 \\
Active months & 0.372 & 0.270 \\

\bottomrule
\end{tabular}
}
\end{table}


AWP exceeds EndorseRank on tenure and active months, as expected from time-decayed transfer edges. Inbound counterparty ratio is weak for both methods. This family is a robustness note, not a headline claim.

\dissertationSubsubheading[sub:trading-success-family]{Trading-success family (not claimed)}

Table~\ref{tab:alignment-gmx} is an artifact table and is not interpreted in the main text.

% Real BigQuery data -- regenerate with:
%   A-Skill-Programs/margin_rank/scripts/run_dissertation_eval.py --real --export-latex
% Generated by export_latex_results.py -- do not edit by hand unless re-export fails
\begin{table}[htbp]
\centering
\caption[Domain alignment with GMX V2 margin-trading success proxies on Arbitrum One.]{Domain alignment with GMX V2 margin-trading success proxies on Arbitrum One. CI = 95\% percentile bootstrap over wallets (400 paired resamples, seed 42).}
\label{tab:alignment-gmx}
\fitwidth{%
\begin{tabular}{lccc}
\toprule
Proxy & Spearman $\rho$ & Kendall $\tau$ & 95\% CI of $\tau$ \\
\midrule
\multicolumn{4}{l}{\textit{EndorseRank}} \\
Close-success count & 0.180 & 0.149 & [0.130, 0.171] \\
Realized gain proxy & 0.132 & 0.107 & [0.086, 0.127] \\
Close success rate & 0.041 & 0.033 & [0.011, 0.055] \\
\addlinespace
\multicolumn{4}{l}{\textit{AWP}} \\
Close-success count & 0.319 & 0.224 & [0.206, 0.241] \\
Realized gain proxy & 0.349 & 0.239 & [0.223, 0.254] \\
Close success rate & 0.110 & 0.075 & [0.055, 0.092] \\

\bottomrule
\end{tabular}
}
\end{table}


\dissertationSubheading[sec:rank-divergence]{Rank divergence (H1 / RQ1)}

Hypothesis H1 and research question RQ1 ask whether EndorseRank and AWP yield meaningfully different wallet orderings on an identical sample. Inter-method Kendall $\tau$ between the two score vectors is $0.316$ on the matched cohort ($n=5{,}521$). The rankings are related but far from interchangeable. Correlation is present because economically active wallets tend to participate in both transfer and approval markets, yet allowance and transfer graphs emphasize different relationships, so many pairwise orderings disagree.

Figure~\ref{fig:score-distributions} shows the marginal distributions of EndorseRank and AWP scores. Both are right-skewed, as is typical for PageRank on sparse directed graphs, but they are not identically shaped. AWP's mass reflects the denser transfer operator; EndorseRank's mass reflects the sparser endorsement operator. Skewness alone does not prove divergence---two methods could share the same ordering with different score scales---but it motivates examining rank-rank relationships directly.

\begin{figure}[htbp]
\centering
\includegraphics[width=0.9\textwidth]{\figpath{score-distributions.pdf}}
\caption{Marginal distributions of EndorseRank and AWP scores on the matched cohort ($n=5{,}521$), shown on a $\log_{10}$ scale for wallets with positive score (EndorseRank excludes $966$ zero-score wallets; AWP excludes $381$). Both are right-skewed PageRank distributions; differences in shape reflect the distinct endorsement and transfer operators.}
\label{fig:score-distributions}
\end{figure}

Figure~\ref{fig:rank-scatter} plots EndorseRank ranks against AWP ranks. Points along the diagonal would indicate identical orderings; the observed cloud shows substantial off-diagonal mass. Wallets that rank highly under transfer centrality are not reliably the same wallets that rank highly under allowance endorsement, and vice versa. The inter-method $\tau=0.316$ summarizes this scatter: agreement exceeds chance, but a protocol that substitutes one ranking for the other would mis-order a large fraction of wallet pairs.

\begin{figure}[htbp]
\centering
\includegraphics[width=0.9\textwidth]{\figpath{rank-scatter.pdf}}
\caption{Rank--rank scatter of EndorseRank versus AWP on the matched cohort. Inter-method Kendall $\tau=0.316$ indicates partial overlap with substantial pairwise disagreement, supporting H1/RQ1 (construct-distinct orderings).}
\label{fig:rank-scatter}
\end{figure}

Rank divergence is the structural counterpart of construct-specific alignment. If the two methods produced nearly identical orderings, family-level winners in Table~\ref{tab:alignment-three-methods} would be difficult to interpret as construct effects; they could be noise around a shared ranking. The observed $\tau=0.316$, together with EndorseRank's allowance lead and AWP's transfer lead, supports reading the methods as related but distinct reputation instruments. H1 is supported: EndorseRank and AWP yield partially overlapping yet construct-distinct orderings on the matched cohort.

An operational reading of $\tau=0.316$ is that for a randomly chosen pair of wallets, the probability that EndorseRank and AWP agree on which wallet ranks higher is only moderately above chance. Protocols that currently use transfer-based scores cannot assume that switching to allowance-based scores---or blending the two---will preserve existing wallet orderings. Migration would re-rank a substantial fraction of pairs, with the largest disagreements likely concentrated among wallets that are transfer hubs but rarely endorsed as spenders, or endorsed spenders with limited transfer footprints. Figure~\ref{fig:rank-scatter} is consistent with that interpretation: the cloud is positively sloped but wide.

The divergence result also disciplines rhetoric about ``on-chain reputation'' as a singular object. On this cohort, reputation-from-authorization and reputation-from-transfer are empirically distinct enough that treating them as interchangeable would misstate the evidence. Chapter~\ref{ch:discussion} returns to the practical implication: method choice should follow the trust construct a protocol needs, not a generic preference for any PageRank variant.

\dissertationSubheading[sec:gf-pr-ceiling]{Temporal holdout of future new approvals (RQ4)}

Scores are frozen at 28~February~2026. Labels are new owner--spender approval pairs first observed from 1~March~2026 through 31~May~2026, counted on the spender. The claimed cohort is t1 inbound spenders ($n=1{,}335$), of whom 222 received at least one new t2 pair.

EndorseRank Kendall $\tau$ with future new approvers is $0.479$ (bootstrap 95\% CI $0.426$--$0.529$; 200 resamples). AWP is $0.044$ (CI $-0.011$--$0.095$). The result is the same endorsement construct in the future tense. It is stronger as out-of-window evidence than the same-window allowance $\tau=0.355$, because the labels are not built from the scored events. It is not a prediction of trading profit, liquidation, or malicious spenders.

The GMX-matched intersection of this label is sparse (12 events) and is not used as a claim. Revoke, drain-heuristic, and Aave-borrower labels were computed and are not reported as support: revoke and drain had the wrong sign for an endorsement reading, and Aave events on this spender set were empty or the wrong construct.

\dissertationSubheading[sec:empirical-claims]{Primary claims (C1--C4)}

Chapter~4 empirical results support the following claims, aligned with Table~\ref{tab:hrq-claim-map}.

\begin{enumerate}
\item[\textbf{C1.}] EndorseRank runtime is lower than AWP runtime on identical $n$. On the matched cohort ($n=5{,}521$), EndorseRank completes PageRank in 2.105\,s versus 18.711\,s for AWP (speedup $\approx 8.9\times$), with 14{,}727 vs.\ 137{,}087 edges (Table~\ref{tab:benchmark-runtime}). On the expanded wallet pool, speedups remain $\approx 10.4\times$ at $n=10{,}000$ (1.18\,s vs.\ 12.28\,s) and $\approx 9.4\times$ at $N=31{,}612$ (2.29\,s vs.\ 21.40\,s) (Table~\ref{tab:benchmark-scaling}).

\item[\textbf{C2.}] EndorseRank and AWP recover their intended constructs on same-window local statistics: EndorseRank allowance mean $\tau=0.355$; AWP transfer mean $\tau=0.534$ (Tables~\ref{tab:alignment-transfer}--\ref{tab:alignment-allowance}). These are construct checks, not trading-skill tests.

\item[\textbf{C3.}] \textbf{Construct-specific alignment:} EndorseRank exceeds AWP on allowance ($\tau=0.355$ vs.\ $-0.030$); AWP exceeds EndorseRank on transfer ($\tau=0.534$ vs.\ $0.333$). Inter-method rank agreement is only moderate ($\tau=0.316$) (Section~\ref{sec:rank-divergence}).

\item[\textbf{C4.}] \textbf{Efficiency--alignment trade-off:} EndorseRank offers $\approx 8.9\times$ lower PageRank runtime than AWP on the matched cohort and $\approx 9$--$10\times$ on expanded-pool scaling while remaining the only social method that aligns with authorization checks. On the spender holdout ($n=1{,}335$), EndorseRank $\tau=0.479$ with future new approvers (CI $0.426$--$0.529$) versus AWP $\tau=0.044$ (Section~\ref{sec:gf-pr-ceiling}).
\end{enumerate}

RQ1 is addressed by inter-method $\tau=0.316$. RQ2 is addressed by allowance versus transfer checks. RQ3 is addressed by the runtime tables. RQ4 is addressed by the spender holdout, not by trading-success families.

What the chapter does \emph{not} claim is equally important. It does not claim that either score predicts trading profit, liquidation, or malicious spenders. It does not claim that AWP is obsolete; AWP remains the stronger method for transfer-hub detection. The contribution is a construct-aware comparison: EndorseRank is cheaper, aligns with allowance checks, diverges moderately from AWP, and predicts later new approvals on the spender holdout.

Chapter~\ref{ch:discussion} interprets these findings against RQ1--RQ4 and discusses implications, threats to validity, and limitations.
